\chapter*{Appendix B: Key Legislation and Regulatory Instruments}
\label{chap:Appendix B: Key Legislation and Regulatory Instruments}

\paragraph{This appendix catalogues the primary legislation, regulations, directives, and regulatory instruments referenced in this book, organised by subject area. It is accurate as of the date of publication but should be read subject to the legislative and regulatory changes that characterise this field.}

%---------------------------- SECTION ----------------------------
\section*{B.1 Health, Safety, and Environment}
\paragraph{\textbf{Health and Safety at Work etc. Act 1974 (HSWA)} - The primary legislation governing occupational health and safety in the United Kingdom. Imposes a general duty on employers to ensure the health, safety, and welfare of their employees so far as is reasonably practicable, and duties to non-employees affected by their work activities.}
\paragraph{\textbf{Management of Health and Safety at Work Regulations 1999 (MHSWR)} - Regulations made under the HSWA requiring employers to conduct suitable and sufficient risk assessments, implement preventive and protective measures, and establish appropriate management arrangements for health and safety.}
\paragraph{\textbf{The Workplace (Health, Safety and Welfare) Regulations 1992} Regulations governing minimum standards for the working environment — covering ventilation, temperature, lighting, cleanliness, space, and welfare facilities.}
\paragraph{\textbf{The Provision and Use of Work Equipment Regulations 1998 (PUWER)} - Regulations governing the safety of work equipment — requiring that equipment is suitable for its purpose, maintained in good condition, and used only by trained and competent persons.}
\paragraph{\textbf{The Manual Handling Operations Regulations 1992} - Regulations requiring employers to avoid the need for manual handling operations that involve risk of injury, and where unavoidable, to reduce the risk so far as is reasonably practicable.}
\paragraph{\textbf{The Control of Substances Hazardous to Health Regulations 2002 (COSHH)} - Regulations requiring employers to prevent or adequately control exposure to substances hazardous to health — encompassing assessment of health risks, implementation of control measures, and health surveillance.}
\paragraph{\textbf{The Regulatory Reform (Fire Safety) Order 2005} - The primary legislation governing fire safety in non-domestic premises — requiring a responsible person to carry out a fire risk assessment and implement appropriate fire safety measures.}
\paragraph{\textbf{The Control of Major Accident Hazards Regulations 2015 (COMAH)} - Regulations implementing the EU Seveso III Directive in Great Britain — requiring operators of sites where significant quantities of dangerous substances are present to prevent major accidents and limit their consequences to people and the environment.}
\paragraph{\textbf{The Environmental Permitting (England and Wales) Regulations 2016} - Regulations establishing the environmental permitting system for activities with significant environmental impact — requiring operators to obtain and comply with environmental permits for regulated activities.}
\paragraph{\textbf{The Environmental Damage (Prevention and Remediation) (England) Regulations 2015} - Regulations implementing the Environmental Liability Directive — imposing operator liability for the prevention and remediation of environmental damage to protected species, habitats, water, and land.}
\paragraph{\textbf{The Construction (Design and Management) Regulations 2015 (CDM)} - Regulations governing health and safety in construction projects — establishing roles and responsibilities for clients, principal designers, principal contractors, designers, and contractors.}
\paragraph{\textbf{The Work at Height Regulations 2005} - Regulations requiring employers and others to plan, supervise, and carry out all work at height in a way that is, so far as reasonably practicable, safe.}

%---------------------------- SECTION ----------------------------
\section*{B.2 Data Protection and Cybersecurity}
\paragraph{\textbf{General Data Protection Regulation (EU) 2016/679 (GDPR)} - The European Union's primary data protection regulation — setting out principles for the lawful processing of personal data, the rights of data subjects, and the obligations of controllers and processors}
\paragraph{\textbf{Data Protection Act 2018 (DPA 2018)} - The primary UK data protection legislation — incorporating the UK GDPR (as retained EU law), implementing the Law Enforcement Directive, and extending data protection requirements to activities not covered by the UK GDPR.}
\paragraph{\textbf{Network and Information Systems (NIS) Regulations 2018} - Regulations implementing the EU NIS Directive in the United Kingdom — imposing security requirements and incident reporting obligations on operators of essential services and relevant digital service providers.}
\paragraph{\textbf{NIS2 Directive (EU) 2022/2555} - The EU directive significantly strengthening and expanding the scope of the original NIS Directive — extending to a wider range of sectors, imposing more stringent security requirements, and enhancing enforcement powers.}
\paragraph{\textbf{Digital Operational Resilience Act (DORA) — Regulation (EU) 2022/2554} - The EU regulation establishing a comprehensive framework for digital operational resilience in the financial sector — covering ICT risk management, incident reporting, resilience testing, and third-party ICT risk management. Applicable from January 2025.}
\paragraph{\textbf{EU AI Act — Regulation (EU) 2024/1689} - The EU's comprehensive regulatory framework for artificial intelligence — classifying AI systems by risk level and imposing proportionate requirements including conformity assessments, transparency obligations, and prohibitions on certain high-risk applications.}

%---------------------------- SECTION ----------------------------
\section*{B.3 Financial Services Regulation}
\paragraph{\textbf{Financial Services and Markets Act 2000 (FSMA 2000)} - The primary legislation governing financial services regulation in the United Kingdom — establishing the regulatory framework, the regulatory objectives, and the powers of the FCA and PRA.}
\paragraph{\textbf{Financial Services Act 2012} - The legislation establishing the Financial Policy Committee, the Prudential Regulation Authority, and the Financial Conduct Authority as successor bodies to the Financial Services Authority.}
\paragraph{\textbf{Financial Services Act 2021} - Legislation making amendments to the UK financial services regulatory framework following Brexit — including changes to the resolution regime and the regulatory framework for benchmarks.}
\paragraph{\textbf{The Money Laundering, Terrorist Financing and Transfer of Funds (Information on the Payer) Regulations 2017 (MLR 2017)} - The primary AML regulatory framework in the United Kingdom — implementing the EU Fourth Money Laundering Directive and imposing customer due diligence, risk assessment, suspicious activity reporting, and governance requirements on regulated firms.}
\paragraph{\textbf{Proceeds of Crime Act 2002 (POCA)} - The primary legislation criminalising money laundering in the United Kingdom and establishing the suspicious activity reporting regime. Creates principal money laundering offences, failure to disclose offences, and tipping off offences.}
\paragraph{\textbf{Terrorism Act 2000} - The primary counter-terrorism legislation — including provisions criminalising terrorist financing and the failure to disclose terrorist financing.}
\paragraph{\textbf{Bribery Act 2010} - The UK legislation creating offences of bribing another person, being bribed, bribing a foreign public official, and — for commercial organisations — failing to prevent bribery. The section 7 corporate offence applies to organisations whose associated persons commit bribery anywhere in the world to benefit the organisation.}
\paragraph{\textbf{Criminal Finances Act 2017} - Legislation introducing corporate offences of failure to prevent the facilitation of tax evasion — both UK tax evasion and foreign tax evasion — and creating unexplained wealth orders.}
\paragraph{\textbf{Market Abuse Regulation (MAR) — Regulation (EU) 596/2014 (as retained in UK law)} - The UK's primary market abuse regulation — prohibiting insider dealing, market manipulation, and the improper disclosure of inside information in relation to financial instruments traded on regulated markets.}
\paragraph{\textbf{Consumer Credit Act 1974 (as amended)} - The primary legislation governing consumer credit in the United Kingdom — regulating credit agreements, hire agreements, and ancillary credit businesses.}
\paragraph{\textbf{Payment Services Regulations 2017} - Regulations implementing the EU Payment Services Directive 2 (PSD2) in the United Kingdom — governing the provision of payment services, the rights of payment service users, and the obligations of payment service providers.}
\paragraph{\textbf{Electronic Money Regulations 2011} - Regulations governing the issuance of electronic money in the United Kingdom — establishing authorisation requirements, capital requirements, and conduct of business obligations for electronic money institutions.}
\paragraph{\textbf{Senior Managers and Certification Regime (SM\&CR)} - Not a single piece of legislation but a regulatory framework implemented through rules in the FCA Handbook and PRA Rulebook under powers in FSMA 2000 — establishing individual accountability requirements for senior managers, certification requirements for significant individuals, and conduct rules for all staff in regulated firms.}

%---------------------------- SECTION ----------------------------
\section*{B.4 Corporate and Commercial Law}
\paragraph{\textbf{Companies Act 2006} - The primary legislation governing companies in the United Kingdom — covering incorporation, share capital, directors' duties, accounts and audit, and insolvency.}
\paragraph{\textbf{Modern Slavery Act 2015} - Legislation consolidating and strengthening existing modern slavery offences — creating offences of slavery, servitude, forced or compulsory labour, and human trafficking — and imposing transparency obligations on large commercial organisations operating in the United Kingdom.}
\paragraph{\textbf{Corporate Manslaughter and Corporate Homicide Act 2007} - Legislation creating the offence of corporate manslaughter — applicable where the way in which an organisation's activities were managed or organised caused a person's death and amounted to a gross breach of a duty of care.}
\paragraph{\textbf{Sanctions and Anti-Money Laundering Act 2018} - Legislation providing powers to make regulations imposing UK sanctions — following Brexit — and amending the anti-money laundering framework.}
\paragraph{\textbf{Electronic Communications Act 2000} - Legislation governing electronic signatures and the admissibility of electronic evidence.}

%---------------------------- SECTION ----------------------------
\section*{B.5 Employment and Equality}
\paragraph{\textbf{Equality Act 2010} - The primary equality legislation in England, Scotland, and Wales — prohibiting discrimination on the basis of protected characteristics including age, disability, gender reassignment, marriage and civil partnership, pregnancy and maternity, race, religion or belief, sex, and sexual orientation.}
\paragraph{\textbf{Employment Rights Act 1996} - The primary legislation governing employment rights in the United Kingdom — including rights relating to unfair dismissal, redundancy, written statements of employment particulars, and protected disclosures (whistleblowing).}
\paragraph{\textbf{Public Interest Disclosure Act 1998} - Legislation protecting workers who make qualifying disclosures in the public interest — providing protection from detriment and dismissal for whistleblowers.}
\paragraph{\textbf{Working Time Regulations 1998} - Regulations implementing the EU Working Time Directive — governing maximum weekly working hours, rest breaks, rest periods, and annual leave entitlements.}

%---------------------------- SECTION ----------------------------
\section*{B.6 Environmental and Sustainability}
\paragraph{\textbf{Climate Change Act 2008} - The primary UK climate change legislation — establishing the legally binding target of net zero greenhouse gas emissions by 2050, the carbon budgeting framework, and the Committee on Climate Change.}
\paragraph{\textbf{Environment Act 2021} - Major environmental legislation — establishing new environmental targets, creating the Office for Environmental Protection, introducing biodiversity net gain requirements, and reforming the environmental governance framework post-Brexit.}
\paragraph{\textbf{Corporate Sustainability Reporting Directive (CSRD) — Directive (EU) 2022/2464} - The EU directive significantly expanding and standardising sustainability reporting obligations for large companies and listed companies — requiring reporting in accordance with the European Sustainability Reporting Standards (ESRS).}
\paragraph{\textbf{Corporate Sustainability Due Diligence Directive (CSDDD) — Directive (EU) 2024/1760} - The EU directive requiring large companies to conduct human rights and environmental due diligence across their operations and value chains.}
\paragraph{\textbf{EU Taxonomy Regulation — Regulation (EU) 2020/852} - The EU regulation establishing a classification system for environmentally sustainable economic activities — providing a framework for identifying which activities contribute substantially to one or more of the six environmental objectives.}

%---------------------------- SECTION ----------------------------
\section*{B.7 Legal Services}
\paragraph{\textbf{Legal Services Act 2007} - The primary legislation governing the regulation of legal services in England and Wales — establishing the Legal Services Board as the oversight regulator, defining the reserved legal activities, and enabling the licensing of alternative business structures.}
\paragraph{\textbf{Solicitors Act 1974} - The primary legislation governing the solicitors' profession — as significantly amended by the Legal Services Act 2007.}